\documentclass{article}

\usepackage{amsmath}
\usepackage{url}
\usepackage{booktabs}
\usepackage{epsfig}
\usepackage{graphics}
\usepackage{latexsym}
\usepackage{amssymb}
\usepackage{comment}
\usepackage{fancyvrb}
\usepackage{moreverb}
\usepackage{mathptmx}
\usepackage[T1]{fontenc}
\usepackage[utf8]{inputenc}
\usepackage{latexsym,color}
\usepackage[french]{babel}
\usepackage{pstricks,pst-node,pst-text,pst-3d,pst-tree,pst-grad}
\usepackage{xcolor}


\newcommand{\fig}[2]
{
 \begin{figure*}[hbt]
  \begin{center}
  \includegraphics[width=10cm]{#1}
  \caption{\label{#1}  #2}
  \end{center}
 \end{figure*}
}

% small et verbatim: to be coherent with \file
\newenvironment{smallverbatim}%
{\endgraf\small\verbatimtab}{\endverbatimtab}


\definecolor{mygray}{gray}{0.90}
\definecolor{mygreen}{rgb}{0,0.5,0}
\definecolor{maroon}{rgb}{0.5,0,0}
\usepackage{listings}


\title{Questions et exercices -- Recherche d'Information}
\author{Philippe Rigaux}

\date{\today}

\begin{document}

\maketitle

\textbf{Indiquez votre nom}\,:


\begin{enumerate}

\section{Généralités}

\item Soit un système qui affiche systématiquement 20 documents, ni plus ni moins, pour toutes les recherches. Indiquez
quel est la précision et le rappel dans les cas suivants\,:

\begin{enumerate}
  \item Je cherche un document unique, il est affiché parmi les 20.
  \item Ma base contient 100 documents, je veux tous les obtenir, le système m'en renvoie 20.
  \item Je fais une recherche dont je sais qu'elle devrait me ramener 30 documents. Parmi les
     20 que renvoie le système, je n'en retrouve que 10 parmi ceux attendus.
\end{enumerate}

\item \textbf{Matrice d'incidence}. Faire un fichier Excel (ou l'équivalent) qiu calcule la
   taille d'une collection, la taille de la matrice d'incidence, et la densité de 1 dans cette
   matrice, en fonction des paramètres suivants\,:
     \begin{enumerate}
       \item $n$, le nombre de documents,
       \item $d$, le nombre de termes distincts,
       \item $m$, le nombre moyen de mots dans un document,
       \item $w$, la taille moyenne d'un mot (en octets).
     \end{enumerate}
     


\section{Recherche avec un index inverse}

\item Essayons de mesurer l'intérêt de trier les listes inversées en comptant le nombre
   de \emph{comparaisons} de docId nécessaires pour évaluer une requête Booléenne.
   
   Prenons pour commencer l'exemple des listes sur \emph{Brutus} et \emph{César} données
    en cours.
    
    \begin{enumerate}
      \item Combien de comparaisons si les listes ne sont pas triées\,? 
      \item Et si les listes sont triées
    \end{enumerate} 
    
    Maintenant, généralisez. Si la première liste a une taille $n$ et la seconde
    une taille $m$, quel est le coût dans les deux cas\,?
    
    \vspace*{4cm}
    
    \item Donnez l'algorithme pour évaluer la négation (``je cherche les pièces avec César mais pas Brutus'').
    
    
    \vspace*{6cm}
    
  \item Le site   \url{http://www.rhymezone.com/} permet de faire des recherches sur les pièces de Shakespeare. 
     Essayez-le et dites comment vous pourriez faire mieux après avoir suivi le cours\,!
     
    \vspace*{4cm}
    
    \item \textbf{Positions}. Voici deux listes inversées avec des positions.
    
    \begin{smallverbatim}
     (be: 993427) 1: [7, 18, 33, 72, 86, 231]; 2: [3, 149]; 4: [17, 191, 291, 430, 434];  …
     (to, 187654) 2: [1,17,74,222,551]; 4: [8,16,190,429,433]; 7: [13,23,191]
    \end{smallverbatim}
    
    \begin{enumerate}
      \item Comprenez-vous la signification des différents chiffres\,? 
      \item Si oui, quel document correspond
     (sans doute) à \textit{Hamlet} (... \textit{to be or not to be}...)\,?
     \item Vous sentez-vous capable de décrire l'algorithme de recherche de phrases
        (tenant compte de la position)\,?
     \end{enumerate}
     
     
\section{Construction d'un index}


\item \textbf{Tri}. On suppose qu'une page contient seulement deux enregistrements et qu'on
   dispose de 4 pages en mémoire. Appliquez l'algorithme de tri-fusion sur le jeu de données 
      suivant (on le lit de gauche à droite puis de haut en bas, et on veut trier sur le numéro).

\begin{smallverbatim}
3 Allier  36 Indre   18 Cher   75 Paris    39 Jura
9 Ariège  81 Tarn    11 Aude   12 Aveyron  25 Doubs 
73 Savoie  55 Meuse   15 Cantal 51 Marne    42 Loire 
40 Landes 14 Calvados 30 Gard   84 Vaucluse 7 Ardèche
\end{smallverbatim}


\item \textbf{Compression}.
La liste inverse d'un terme $t$ contient les docId suivants, codés sur 4 octets\,:

\[[345; 476; 698; 703].\]

Appliquer la compression \textit{VByte} sur cette séquence. Quel est le gain en espace\,?

\end{enumerate}


\end{document}